\documentclass[11pt]{article}
\usepackage{amsmath,amssymb,amsthm,hyperref}
\newtheorem{theorem}{Theorem}\newtheorem{proposition}[theorem]{Proposition}\newtheorem{definition}[theorem]{Definition}
\title{Scoped control comprehension, dependent grafting, and residual disclosure}
\date{16 November 2026}
\begin{document}\maketitle
\paragraph{Audit and scope.} The written arguments have passed independent review at the declared scope; no proof-assistant certification is claimed. The source is the original finite normalized native context category for the rig $\mathbb R_{\geq0}^{\rm disc}\times\mathcal A$, with cocartesian addition in $\mathcal A$, and $\mathcal E=\mathcal A/I$ possessing finite products and its inherited symmetric tensor. All actual augmentation-preserving coefficient maps are retained. The tensor unit is terminal in $\mathcal E$. Three separate interfaces are proved below: a coupling-control comprehension category; a finite retained-history dependent grafting calculus with coherent directed coefficient evaluation; and disclosure of an already sampled finite scope. None is silently identified with an unrestricted distributive law on unscoped processes. Full original content grades and their supplied maps are separate data; the constructions never replace a grade by its realized support.

\section{A failure before the associator}
At an outer row let $W=\{0,1\}$ be fair, and let $r:W\to\{*\}$ be its native collapse, with unit coefficients. After either outer witness the continuation consists of two fair-bit primitive rows. Let $q_0$ be their diagonal coupling and $q_1$ their antidiagonal coupling. The continuation primitive comparison is identity. A pointwise transport of the selected continuation controls to a policy on the target singleton would require
\[
 q_*=q_0=q_1,
\]
which is impossible. Nevertheless the flattened law, after erasing the outer bit, is the independent fair coupling $(q_0+q_1)/2$.

\begin{proposition}[Section descent obstruction]
The category of coupling controls at one supplied input is functorial under native row comparisons, but the operation of selecting a dependent section after a sampled outer witness does not admit pointwise pushforward along all original native outer maps. This failure occurs for fixed arity two, terminal coefficients, and finite normalized laws. Hence a skew associator, without a changed control interface or a restricted class of cells, cannot repair it.
\end{proposition}
\begin{proof}
The displayed collapse is an original native context map. A section is an actual assignment at every raw witness, including witnesses of weight zero. A pointwise section comparison requires its control equation at both raw points; the two distinct displayed laws contradict that equation. The valid mixture is a different, averaged control equation, not the missing pointwise equation.\end{proof}

If a continuation choice collection contains only $\delta_0,\delta_1$, choosing $\delta_w$ before the same collapse also gives an object-existence failure: the required fair mixture is absent from that collection. Convex completion repairs that particular existence failure, but does not identify actual raw witnesses, coefficient arrows, or zero-history conditional choices. In the all-coupling example the mixture is already admitted, yet pointwise descent still fails. These are distinct obstructions. The independent-product policy is constant under this collapse and does not exhibit either failure.

\section{An explicit scoped control category}
Let $\mathscr D$ be a specified category of finite occurrence-row diagrams. A diagram $D$ has visible input set $X_D$, and at $x$ a finite ordered family of normalized native rows. A map $r:D\to D'$ has a visible map $r_X$, a specified ordered-position correspondence (identity for original constant-polynomial shape maps), and actual native row maps. Assume the resulting coupling pushforward is the functorial map
\[
 C(r):C(D)\to C(D'),\qquad C(D)=\coprod_{x\in X_D}\mathsf{Coup}(D_x),
\]
lying over $r_X$. These assumptions hold for the finite constant-polynomial controlled crossing; no general display-monad lift is inferred.

\begin{definition}
A scoped assignment over $D$ is a set $M$ and map $a:M\to C(D)$. Its visible input is the composite $p:M\to X_D$. A morphism $(D,M,a)\to(D',M',a')$ is $(r,k)$ with $k:M\to M'$ and $a'k=C(r)a$. All coefficient maps remain in $r$, not merely in $C(r)$. Denote this category by $\mathsf{Scope}(C)$.
\end{definition}
\begin{theorem}[Split transport and exact descent]
The projection $\mathsf{Scope}(C)\to\mathscr D$ is a split opfibration. Its chosen lift along $r$ sends $(D,M,a)$ to $(D',M,C(r)a)$ and keeps memory $M$ unchanged. If ordinary policies are identified with the special objects $M=X_D$, $p=1$, the pushed scoped assignment descends to an ordinary policy on $X_{D'}$ precisely when $C(r)a$ is constant on every fibre of $r_X$ and choices have been supplied at points outside its image. The descended choice, when it exists, is unique on the image, but not outside it.
\end{theorem}
\begin{proof}
Composition is $(s,l)(r,k)=(sr,lk)$ and its control equation follows by functoriality of $C$. A morphism $(sr,k)$ from $(D,M,a)$ factors uniquely through the chosen lift by $(s,k)$, since its equation is exactly $a''k=C(s)C(r)a$. Identity and composite lifts agree literally. Descent is the elementary factorization of the set map $C(r)a$ through $r_X$, with the prescribed visible base; this factorization holds exactly under the stated fibre equality and extension data.\end{proof}
This is a genuine universal transport property, but its codomain remembers the old memory. It does not assert that erased observations can be recovered from the new visible input.

\subsection{The necessary sampled-memory dilation}
For a single native row map
\[
 r:(W,\mu,C_w)\longrightarrow(W',\mu',C'_{w'}),
 \qquad \alpha_w:C_w\to C'_{r(w)},
\]
define its retained implementation to have raw set $W$, law $\mu$, coefficient $C'_{r(w)}$, and the target body evaluated at $r(w)$. There are native maps
\[
 (W,\mu,C_w)\xrightarrow{(1,\alpha)}
 (W,\mu,C'_{r(w)})\xrightarrow{(r,1)}
 (W',\mu',C'_{w'}).
\]
At context-row level the first changes coefficients, not the retained sample; the second forgets the latent sample. For processes, assume the original comparison already supplies its body square $b^{\prime}r=g b$ for the specified endpoint map $g$. The retained implementation has body $b^{\prime}r$, so the first factor covers $g$ and the second covers the identity target endpoint. No arbitrary change of observable body is permitted; for fixed endpoints $g=1$. Its scalar equation is $r_*\mu=\mu'$, and its actual coefficient arrows are identities. In the original rig the fibre aggregation arrow is the coproduct fold, so this is an actual native factorization, not just a probability equality. It includes every zero point.

For composable $r,s$, the twice-retained final row and the directly retained row both have raw $W$, law $\mu$, and coefficient $C''_{sr(w)}$; their maps agree by composition of the supplied coefficient arrows. If graph tuples are used instead of $W$, their canonical tuple-erasure isomorphisms obey the pentagon because each retains the same original coordinate. This proves coherent finite iterated dilation. It deliberately changes the target implementation: a future controller may inspect latent $w$, whereas a controller on $W'$ alone cannot. Thus the factorization is an explicit way to retain scope, not a proof that arbitrary policies descend.

\section{Retained-history trees and dependent grafting}
A finite scoped law tree has a root, a fixed finite maximum depth, and at every nonterminal raw node $h$ a finite nonempty successor set $W_h$, normalized law $q_h$, and coefficient $K_{h,w}\in\mathcal E$ on each outgoing edge. Every successor, including every zero-weight successor, remains a node. An empty-duration tree consists of its root as a leaf. Each leaf may have an output label; each nonterminal node may carry an original polynomial shape and its ordered active-occurrence data. In the frontier application $W_h$ is the product of the currently active primitive witness sets, $q_h$ is a supplied coupling with their specified conditional marginals, and $K_{h,w}$ is the product of their actual coefficients. All those active inputs are known at $h$, before sampling $w$.

A native tree map is a depth-preserving rooted map which preserves and reflects terminality, whose raw successor maps push each $q_h$ to the target row at the image of $h$, and which supplies each actual edge coefficient arrow. It also preserves the supplied body/shape equations. A shorter terminal branch stays terminal: no probabilistic idle rounds are inserted. The identity and composition are immediate rowwise. Maps with incompatible stopping strata must first be compared at another explicitly specified tree interface; they are not included by assertion.

\begin{theorem}[Dependent grafting, cuts, and coherence]
Let $T$ be such a tree and, for every raw leaf $h$ of $T$, let $S_h$ be another such tree. Grafting forms a scoped tree $T[S]$ by identifying the root of $S_h$ with $h$. Its raw leaves are precisely pairs $(h,z)$ with $z$ a leaf of $S_h$. It retains all the prefix and suffix controls, including those at zero-weight prefixes. This operation is functorial for a native map $f:T\to T'$ together with native maps $S_h\to S'_{f(h)}$. Empty-duration trees give units. Iterated grafting has the canonical associator
\[
 (T[S])[R]\ \cong\ T[\,S_h[R_{h,-}]\,],
\]
which preserves all controls and actual coefficient maps. The associator and units satisfy the pentagon and triangle. Conversely, cutting at a specified raw stopping frontier determines the prefix tree and the entire family $S_h$ uniquely. Here a raw stopping frontier is a complete antichain of nodes meeting every raw root-to-leaf path exactly once, including all zero-weight paths; the root and terminal leaves are allowed.
\end{theorem}
\begin{proof}
Nodes are addressed paths. Grafting concatenates addresses and retains the original transition row and coefficient on every edge. A native map is obtained by the specified prefix and suffix maps; all its conditional equations are inherited from the corresponding edge. Both iterated graftings have the same addressed primitive edges and controls, with only the brackets in their addresses changed. All five pentagon paths send a nested address to the same concatenated address and act identically on primitive edge maps. The triangle is the empty-address identity. A stopping frontier partitions all nodes into the prefix and the descendants of its raw frontier points, including zero points; this partition proves the inverse cut statement.\end{proof}

This theorem specifies an indexed, dependent operation: a suffix is supplied for every \emph{raw leaf}, not merely for every visible output label. It is not a binary operation on ordinary unscoped input/output processes. If suffix data are supplied only at output labels, they can be reindexed to leaves; the section-descent proposition explains why an arbitrary leaf-indexed family cannot be pushed back to a coarser output interface. Thus the construction supplies the missing comprehension interface rather than calling an undefined operation a skew category.

\subsection{Evaluation and directed coefficient laws}
A leaf path $h=(w_1,\ldots,w_t)$ has mass and staged coefficient
\[
 P_T(h)=\prod_{j=1}^t q_{h_{<j}}(w_j),\qquad
 K_T(h)=\bigotimes_{j=1}^t K_{h_{<j},w_j}.
\]
The empty tensor is $I$. Finite induction proves normalization, and proves exact native leaf pushforward for every native tree map. For grafting, $P_{T[S]}(h,z)=P_T(h)P_{S_h}(z)$ and $K_{T[S]}(h,z)\cong K_T(h)\otimes K_{S_h}(z)$; actual arrow coherence follows from monoidal coherence. Consequently this evaluator respects the preceding dependent grafting, not a nonexistent unscoped policy pushforward.

Suppose an executed frontier has coefficient $K_j=\prod_{i\in I_j}C_{j,i}$. There is a natural arrow
\[
 \bigotimes_j\prod_{i\in I_j}C_{j,i}\longrightarrow
 \prod_{(j,i)}C_{j,i}.\tag{A}
\]
Its $(j,i)$-projection augments every other temporal factor to $I$ and projects the $j$th factor. Finite-product uniqueness proves that arbitrary temporal cuts and frontier-product regroupings commute with (A), including empty frontiers. It need not be invertible. This is the exact bridge to the all-executed-node product coefficient of the causal frontier evaluator.

For fixed occurrence positions $i\in I$ continuing through $t$ stages, the more precise temporal interchange is
\[
 \bigotimes_{j=1}^t\prod_{i\in I} C_{j,i}
 \longrightarrow \prod_{i\in I}\bigotimes_{j=1}^t C_{j,i}.\tag{B}
\]
Each target projection tensors the $i$th source projections. Naturality, associativity for arbitrary temporal blocks, and compatibility with regrouping the fixed occurrence family follow after each target projection from monoidal coherence. At $t=0$ both sides are canonically terminal; at $|I|=1$ the map is the canonical identity. The empty-occurrence case uses the unique map to the terminal object. For $t=1$ it is the identity. These are actual unital lax interchange equations, not a claimed invertible distributive law. The complete controlled probability law is evaluated by the same retained conditional tables under every cut; the exact causal factorization criterion characterizes the image among flat whole-trajectory couplings. Zero-history choices remain in the source and can be forgotten by evaluation.

Equation (B) presupposes a fixed correspondence of a continuing occurrence across the temporal blocks. When variable-arity execution creates new occurrences, (A) remains defined, but (B) is not silently applied to nonexistent corresponding positions. At a supplied complete current shape, grouping occurrence products remains exact; grouping a whole strategy may require the full common history because a group marginal can depend on another group's past. These are the remaining indexed bases, not cosmetic bracketing choices.

\paragraph{Independent locus.} Setting every $q_h$ to the product of its primitive conditional rows gives the original independent scalar bind. Every original pointwise native primitive comparison compatible with the specified routing and stopping interface satisfies the tree-map conditional equations by functoriality of products of finite laws. Hence all those original arrows survive, including noninjective witness maps and nonfaithful coefficient maps. The native coefficient comparison from the original tensor of primitive coefficients to the staged/product evaluation is the augmentation/projection map, and may have a nonzero kernel. The retention claim here is bookkeeping: each supplied original raw, coefficient, and global-grade arrow remains named in the input data. No category of graded scoped trees with all original global-grade arrows and no source inclusion functor has yet been constructed here, so no categorical faithfulness, full faithfulness, or conservativity theorem is claimed. Constructing that graded source inclusion and proving its faithfulness is a separate obligation. The evaluated target has extra shared controls and may forget tensor directions under evaluation.

\section{Residual disclosure of one finite sampled scope}
Fix one supplied finite occurrence family, with all local inputs known before the draw, a joint law $q$ on $W=\prod_{i\in I}W_i$, and native coefficient $K_w=\prod_i C_{i,w_i}$ (or any other explicitly supplied coefficient $K_w$). For each $J\subseteq I$ define the actual retained-state set
\[
 R_J=\{(w,v):w\in W,\ v=w|_J\}.
\]
The initial native context row has raw $W$, law $q$, coefficient $K_w$, and body $w\mapsto(w,w|_J)\in R_J$. Its visible observation is the second projection; the full latent coordinate $w$ remains in the interface $R_J$. This interface must not be replaced by the smaller set of observations.

\begin{theorem}[Residual reveal calculus]
For $J\subseteq J'$, reveal is the pure function $R_J\to R_{J'}$ sending $(w,w|_J)$ to $(w,w|_{J'})$. As a native deterministic context it has one raw witness and unit coefficient at each retained state. Composing the initial sampled row with it retains raw $W$, law $q$, and coefficient $K_w\otimes I\cong K_w$. These reveals compose strictly by subset inclusion. Revealing disjoint blocks in either order gives the same full record and observation, and arbitrary finite disclosure cuts agree. Each original componentwise native scope map $r=\prod_i r_i$ induces a native map of residual records, commuting with every reveal and carrying the actual supplied product coefficient arrow. No additional copy of $K_w$ is created by a reveal.
\end{theorem}
\begin{proof}
Under the canonical identifications $R_J\cong W$, every asserted state equation is identity on the retained sample. The deterministic native rows compose with unit coefficients. On observations the commuting equation uses restriction from the larger disclosed subset to the smaller subset. No inverse extension map on visible observations alone is asserted. The transported law is $r_*q$ and the coefficient square is the original native square; hence transport and reveals commute, not only after passing to laws.\end{proof}
Operationally the initial draw contributes $K_w$ exactly once; all later deterministic reveal operations have unit coefficient. Tensoring these unit factors preserves $K_w$. Discard of latent coordinates is the original raw projection with its supplied actual coefficient map, if such a map is specified; observing only a subset does not itself authorize replacing $K_w$ by a smaller product. A canonical projection $K_w\to\prod_{i\in J}C_{i,w_i}$ exists for the displayed product doctrine, but is directed and can discard coefficient information.

This is the standard retained-state/pre-sampling mechanism, not a new unrestricted interference theorem; the project already has lazy/eager cache precedents. In a causal frontier evaluator it attaches immediately after a particular frontier is sampled: store its whole raw tuple and disclose outputs through unit-cost deterministic steps. The next new random frontier still requires its conditional rows at the actual retained history. The residual theorem does not authorize replacing a fresh conditional draw by a reused unconditional marginal. For example a diagonal fair-bit scope, after its first coordinate is observed, has a deterministic second coordinate, not a fresh fair-bit row.

\paragraph{Boundaries of pre-sampling.} Unknown future inputs are not part of a currently supplied scope. A finite-horizon contingent plan could sample prospective rows for all possible histories, but then off-path raw coordinates and coefficients are additional data; discarding them is generally directed, not a native isomorphism. No such plan compiler is proved here. Infinite-horizon plans may require an uncountable seed space and infinite coefficient products, outside the stated finite native doctrine.

\section{What has and has not assembled}
The coupling bundle admits split scoped transport; retained-history trees admit an explicit dependent substitution with unit and pentagon; temporal coefficient evaluation obeys actual unital lax interchange; residual reveal admits strict disclosure/cut laws. The section-descent counterexample proves why these facts do not imply one crossing on ordinary unscoped processes with all original native maps. A completion must specify whether it retains memory/dilations, uses an averaged comparison and proves coefficient-compatible disintegration, or restricts native policy cells by conditional lumpability. Those choices have different kernels. Merely calling the result lax or skew does not choose between them.

Full global content grades in the recursive application are the supplied marked tree grades before cut erasure. The finite grafting theorem does not identify marked and unmarked grades; a map forgetting cut markers is generally noninvertible. The existing causal frontier theorem supplies the terminating, potentially unbounded evaluator under its atomic-frontier and full-history assumptions. The present finite calculus attaches to each bounded prefix and its scoped cuts; it does not by itself extend any infinite coefficient construction or prove scheduler invariance.
\end{document}
